\documentclass[11pt, a4paper]{article}
\usepackage{amsthm}
\usepackage{amsmath}
\usepackage{amssymb}
\usepackage{geometry}
\geometry{
    margin=1in,
    headheight=20pt,
    footskip=10mm
}
\title{8}
\author{Nagusa, G}
\date{\today}

\begin{document}
\maketitle

\section*{Problem}
Let $C$ be the set of all continuous functions from $\mathbb{R} \to \mathbb{R}$. Define two binary operations
\begin{align*}
    \quad \oplus : C \times C &\to C & \odot : C \times C &\to C \\
    (f, g) &\mapsto f \oplus g & (f, g) &\mapsto f \odot g
\end{align*}
where
\begin{align*}
    (f \oplus g)(x) &:= f(x) + g(x) & (f \odot g)(x) &:= f(x) \cdot g(x) \quad (x\in \mathbb{R})
\end{align*}
It is well known that $<C,\oplus,\odot>$ is a ring\\
Show that
\begin{enumerate}
    \item The subset $S = \{f \in C \mid f(3)=0\}$ is a subring of $C$.
    \item $C$ is not an integeral domain
\end{enumerate}

\section*{Proof}

\subsection*{1.}
Let $S = \{f \in C \mid f(3)=0\}$ be a subset of $C$.\\
We will verify the following to prove that $S$ is a subring of $C$
\begin{enumerate}
    \item $0_C \in S$\\
    Let $0_C$ be the zero function defined by $0_C(x) = 0$ for all $x \in \mathbb{R}$.\\
    $0_C$ is continuous and $0_C(3) = 0$.\\
    Thus, $0_C \in S$.\\
    Moreover, since $0_C(x) + f(x) = f(x)$ for all $f \in S$, $0_C$ is the identity element of $C$.
    
    \item $f \oplus (-g) \in S$ for all $f,g \in S$\\
    Let $f, g \in S$. Then $f(3) = 0$ and $g(3) = 0$. 
    We observe that
    $$(f \oplus (-g))(3) = f(3) + (-g)(3) = f(3) - g(3) = 0 - 0 = 0$$
    Since the difference of two continuous functions is continuous, $f \oplus (-g) \in S$.
    
    \item $f \odot g \in S$ for all $f,g \in S$\\
    Let $f, g \in S$. Then $f(3) = 0$ and $g(3) = 0$.
    We observe that
    $$(f \odot g)(3) = f(3)\cdot g(3) = 0 \cdot 0 = 0$$
    Since the product of two continuous functions is continuous, $f \odot g \in S$.
\end{enumerate}
Thus $S$ is a subring of $C$.

\subsection*{2.}
To prove that $C$ is not an integral domain, we will find zero divisors in $C$.\\
Define $f, g$ as follows
\begin{align*}
    f(x) &= 
    \begin{cases} 
        x & \text{if } x \geq 0 \\ 
        0 & \text{if } x < 0
    \end{cases} \\
    g(x) &= 
    \begin{cases} 
        0 & \text{if } x \geq 0 \\ 
        x & \text{if } x < 0 
    \end{cases}
\end{align*}
Clearly, both $f$ and $g$ are continuous functions on $\mathbb{R}$, so $f, g \in C$. \\
We observe that for any $x \in \mathbb{R}$,
\[
(f \cdot g)(x) =
\begin{cases} 
    x \cdot 0 = 0 & \text{if } x \geq 0 \\ 
    0 \cdot x = 0 & \text{if } x < 0
\end{cases}
\]
Since $(f \cdot g)(x) = 0$ for all $x$, this implies $f \odot g = 0_C$. \\
Because $f \neq 0_C$ and $g \neq 0_C$, both $f$ and $g$ are zero divisors.\\
Thus $C$ is not an integral domain.

\end{document}